\section{Schopenhauer on Christianity}

\begin{quotationx}
The myth of the Fall of man is the only thing in the Old Testament to which I can concede a metaphysical, although only allegorical, truth: indeed it is this alone that reconciles me to the Old Testament. Thus our existence resembles nothing but the consequence of a false step and a guilty desire. New Testament Christianity, the ethical spirit of which is that of Brahmanism and Buddhism, and which is therefore very foreign to the otherwise optimistic spirit of the Old Testament, has also, extremely wisely, started from that very myth ... The innermost kernel and spirit of Christianity is identical with that of Brahmanism and Buddhism.

The spirit and ethical tendency, however, are the essentials of a religion, not the myths in which it clothes them.

Therefore, I do not abandon the belief that the teachings of Christianity are to be derived in some way from those first and original religions i.e., Brahmanism and Buddhism . ... But in virtue of this origin, Christianity belongs to the ancient, true, and sublime faith of mankind. This faith stands in contrast to the false, shallow, and pernicious optimism that manifests itself in Greek paganism, Judaism, and Islam.


\flright{\textsc{Arthur Schopenhauer}}
\end{quotationx}


Schopenhauer was a student of the \emph{Upanishads}, since he found those works --- along with those of Buddhism --- to be compatible with his philosophical outlook. What may be less well known, is that he put Christian mysticism, at least that based on a Neoplatonic outlook which was very influential in the early church, in the same class. Unlike the many \emph{cultured despisers of religion}, Schopenhauer looked past the dogmas, and instead focused on the mental and spiritual attitudes engendered by a religious system.

\url{https://www.youtube.com/watch?v=kLyRoKzUCaA} 

\flrightit{Posted on 2009-08-06 by Cologero}

\begin{center}* * *\end{center}

\begin{footnotesize}
\begin{sffamily}
\texttt{OMAR (@OAS94) on 2019-09-10 at 23:14 said:}

Completely taken out of context. It is well known that Schopenhauer despised Christianity (Abrahamic religions in general)

\hfill

\texttt{OMAR (@OAS94) on 2019-09-10 at 23:18 said:}

``But if any one should believe for this reason that European morals have improved proportionally, and that now at any rate they surpass what obtains elsewhere, it would not be difficult to demonstrate that among the Mohammedans, Gnebres, Hindus, and Buddhists, there is at least as much honesty, fidelity, toleration, gentleness, beneficence, nobleness, and self-denial as among Christian peoples.''

On the Basis of Morals, Schopenhauer

\hfill

\texttt{OMAR (@OAS94) on 2019-09-10 at 23:20 said:}

``Indeed, the scale will be found rather to turn unfavourably for Christendom, when we put into the balance the long list of inhuman cruelties which have constantly been perpetrated within its limits and often in its name. We need only recall for a moment the numerous religious wars; the crusades that nothing can justify; the extirpation of a large part of the American aborigines, and the peopling of that continent by negroes, brought over from Africa, without the shadow of a right, torn from their families, their country, their hemisphere, and, as slaves, condemned for life to forced labour; 2 the tireless persecution of heretics; the unspeakable atrocities of the Inquisition, that cried aloud to heaven; the Massacre of St. Bartholomew; the execution of 18,000 persons in the Netherlands by the Duke of Alva; and these are but a few facts among many.''

On the basis of morals

\hfill

\texttt{Cologero on 2019-09-10 at 23:35 said:}

Gee whiz, Omar. You need to stay away from the Bronze Age Pervert. What do you bench?

Intelligent readers will recognize that the quote was entirely in context. The teachings cannot be judged by the crimes of its nominal adherents.

\hfill

\texttt{Tom on 2021-12-21 at 22:07 said:}

Who is the '' Bronze Age Pervert'' ?

\end{sffamily}
\end{footnotesize}

